% !TEX root = ../../../main.tex
% Block card: preprocessing.invert -- "Invert" (Part I, the blocks chapter).
\cardfield{What it does}
Turns the signal upside down: every value is multiplied by $-1$.

\cardfield{When to use it}
To reuse a detector built for one direction on events that point the other way. Put it before
the drop detector and every spike becomes a drop, so the drop detector finds spikes, much as
you might turn a photo upside down to check the shape of a valley as if it were a hill. Then
tell \blockref{interrogation.event_shape} that its input was inverted (its
\emph{upstream\_inverted} setting), so it reports the events as the spikes they really are.

\cardfield{What goes in, what comes out}
In goes a span of one channel; out comes the same span flipped, same length and rate.

\cardfield{Settings that matter}
None.

\cardfield{Watch out for}
Anything measured downstream is in the flipped frame: a depth is a height, a fall is a rise.
Only \blockref{interrogation.event_shape} knows how to flip back, and only when you tell it.
